%=============================================================================
\section{Implications for anticipatory governance and risk mitigation}
\label{sec:implications}
%=============================================================================

The results support a specific role for BVS in the governance of a mandated
transition and exclude others. The system reproduces what national data
already measure, and it organises thin external evidence into explicit,
queryable scenarios for what they do not. It does not measure Spanish
acceptance. Its outputs are therefore inputs to risk mitigation, used to
decide where to pilot, whom to engage, and what to measure first. This section
sets out that use. The governance contribution demonstrated here is the
organisation of evidence and assumptions into a decision-support workflow;
improvements in decision quality, user understanding, and implementation
outcomes remain to be evaluated.

\subsection{From scenarios to a risk-mitigation procedure}

Following the distinctions of Section~\ref{sec:theory}, mitigation proceeds as
a cycle and not as a single prediction-to-policy step:
\emph{explore} scenarios, \emph{identify} the assumptions that drive them,
\emph{deliberate} with affected actors, \emph{pilot} flanking measures,
\emph{measure} the target outcomes, and \emph{update} the model.
Table~\ref{tab:mitigation} maps the main adoption risks of the
recycled-content mandate onto this cycle.

\begin{table}[htbp]
\centering
\caption{Risk-mitigation matrix for recycled-content food-contact packaging:
what BVS contributes, which measure it suggests testing, and which evidence is
needed before acting at scale.}
\label{tab:mitigation}
\small
\begin{tabularx}{\textwidth}{>{\raggedright\arraybackslash}p{2.9cm}XXX}
\toprule
\textbf{Risk} & \textbf{Contribution of BVS} & \textbf{Measure to test} &
  \textbf{Evidence required} \\
\midrule
Low acceptance in some segments & Indicates candidate segments under stated
  assumptions & Communication and product-experience pilots in those segments
  & Target-specific surveys; observed purchases \\
Distrust of safety & Flags safety perception as a priority for measurement &
  Verifiable information, traceability, certification marks, comprehension
  tests & Communication experiments; technical safety evidence under
  Regulation~2022/1616 \\
Wrong transfer of foreign evidence & Shows sensitivity to base rate and
  transfer factor & Collect Spanish data on the target outcomes & Dedicated
  validation of acceptance \\
Unequal burden across groups & Disaggregates scenarios and intervals by
  profile & Distributional review with consumer organisations & Impact
  assessment by group \\
Over-confident use of the interface & Displays assumptions, sources and
  validation status with each answer & Human oversight, query logging, option
  to return no estimate & Evaluation of parsing and interpretation errors \\
\bottomrule
\end{tabularx}
\end{table}

\paragraph{Targeting flanking measures.}
A spread of several tens of points between profiles, if confirmed, means that
one communication strategy cannot serve all segments. The scenarios justify
comparing instruments in pilots before rollout: safety-framed information and
third-party certification against environmental-benefit messaging, and
traceability (for instance, access to recycling-process documentation)
against price measures. They do not establish which instrument works. The
design identifies no causal effect of any intervention, and stated acceptance
is not purchase.

\paragraph{Staging within the legal timeline.}
The surface can inform the order in which firms introduce recycled content
across product lines and channels, starting where the purchasing segment
resembles receptive profiles and using the experience gained to prepare
categories closer to Scenario~B. Staging concerns commercial sequencing and
pilots ahead of the 2030 targets. It is no argument for postponing legal
obligations.

\paragraph{Prioritising measurement.}
The most direct implication is where to spend fieldwork. The held-out tests
show that levels do not transfer, and the target outcomes rest on one or two
sources each. A Spanish survey module on food-contact acceptance, perceived
safety, and willingness to pay would replace the weakest assumptions of the
model, and Eq.~\eqref{eq:fusion} absorbs it without rebuilding. BVS indicates
which segments such a module should oversample. Table~\ref{tab:plan} turns this
into a sampling plan: for each segment, the number of Spanish respondents
needed for a 95\% posterior interval of $\pm 5$ points. Under the
transport-based prior the transferred evidence is worth about 15--25
respondents per segment, so the first Spanish measurement carries almost all of
the information: between about 100 and 370 respondents per segment, so a module
of 1{,}500 respondents with oversampling of higher-educated, non-recycling, and
unworried respondents covers all of them.
Pilots follow the same logic: an intervention is compared between segments of
contrasting projected acceptance, each sized from the table, and the result
enters the model through Eq.~\eqref{eq:fusion}.

\begin{table}[htbp]
\centering
\caption{Projected food-packaging acceptance (\%, 95\% interval) and
measurement plan by segment. \emph{Parameter}: uncertainty of the published
coefficients and base rate only. \emph{Transport}: adds the empirical
transport error of Section~\ref{sec:transport}. $m$, $m_T$: Beta concentration
matched to each set of draws. $N$, $N_T$: Spanish respondents needed for a
95\% posterior interval of half-width 5 points under each prior. Sources:
\texttt{transfer\_output/measurement\_plan.csv} and
\texttt{transport\_segments.csv}.}
\label{tab:plan}
\small
\setlength{\tabcolsep}{4pt}
\begin{tabular}{lrrrrrr}
\toprule
\textbf{Segment} & \textbf{Parameter} & \textbf{Transport} & $m$ & $m_T$ & $N$ & $N_T$ \\
\midrule
%%PLANROWS%%
\bottomrule
\end{tabular}
\end{table}

\subsection{Users, decisions, and responsibilities}

Three groups of users have distinct needs. \emph{Regulators and public
agencies} can use the scenarios to plan information campaigns and to decide
which outcomes national statistics should start to measure. \emph{Firms},
including the small and medium enterprises that make up much of the
food-contact packaging sector and cannot commission repeated fieldwork, can use
them to choose where to pilot. \emph{Consumer and civil-society organisations}
can use the disaggregated scenarios to question whose acceptance is being
assumed. For all three, the interface lowers the cost of access (R4), which is
also what makes misuse easy.

\paragraph{How the system can shape and bias decisions.}
The outputs bear on four decisions: which segments are prioritised, how a
measurement budget is allocated, which pilots are run and where, and which
outcomes national statistics begin to track. Each can be distorted. Segments
absent from the source study are ranked only through the transferred
coefficients; omitted predictors make the surface look flatter than it may be;
a measurement budget steered by the model is spent where the model already
expects differences; and a probability reported to the percentage point lends
weak assumptions an authority they do not have. Reporting the sensitivity range
with every estimate, reserving part of the measurement budget for segments the
model does not flag, and stating the source population with each answer limit
these effects.

\subsection{Governing the instrument}

Used in or around rule-making, BVS is a technology with its own governance
requirements \citep{stilgoe2013}. Four conditions follow from the results.

\emph{Provenance.} Every estimate must be traceable to sources, extracted
statistics, transformations, and model version. The pipeline exports these
tables and fixes its random seeds; the audit of the current evidence base
(Section~\ref{sec:limitations}) shows why the tables must be reviewed and not
only produced.

\emph{Separation of roles.} Language models extract and translate; they must
not be the origin of a number. Extraction errors propagate silently into
priors, so extraction requires human sign-off.

\emph{Status labelling.} Each answer should state whether the outcome is
observed in national data, validated externally, or projected from transferred
evidence, and the system should decline queries outside its domain.

\emph{Contestability and accountability.} Affected parties need a route to
challenge sources and mappings, and the body that acts on an estimate remains
answerable for the decision. Synthetic respondents do not consent or
deliberate. Their use must not displace consultation, and a projection must
not be cited as evidence that the public accepts a measure.

\subsection{Beyond recycled plastics}

Mandates that precede behavioural measurement are common: novel food
ingredients, alternative proteins, reuse-and-refill systems, extended producer
responsibility. The BVS architecture requires only a national survey with
demographic and attitudinal variables, external studies reporting aggregate
outcomes for the target constructs, and a defensible ordering among variable
blocks. Unlike approaches that improve as early market data accumulate
\citep{lee2012}, it is built for the window before such data exist. Its
transfer to other domains is an architectural possibility. Its performance
there has not been evaluated, and the transfer tests reported here would need
to be repeated in each case.

